% Generated draft: requires mathematical and source review.
\documentclass[11pt]{article}
\usepackage[margin=25mm]{geometry}
\usepackage{fontspec}
\setmainfont{Latin Modern Roman}
\usepackage{amsmath,amssymb,mathtools,unicode-math}
\setmathfont{Latin Modern Math}
\usepackage{enumitem,needspace,xurl,hyperref}
\setlength{\parindent}{0pt}
\setlength{\parskip}{5pt}
\setlength{\emergencystretch}{3em}
\newcommand{\N}{\mathbb{N}}
\newcommand{\Z}{\mathbb{Z}}
\newcommand{\Q}{\mathbb{Q}}
\newcommand{\R}{\mathbb{R}}
\newcommand{\C}{\mathbb{C}}
\newcommand{\F}{\mathbb{F}}
\newcommand{\A}{\mathbb{A}}
\newcommand{\PP}{\mathbb{P}}
\newcommand{\E}{\mathbb{E}}
\newcommand{\eps}{\varepsilon}
\newcommand{\dd}{\,\mathrm{d}}
\begin{document}
\section*{UnsolvedMath ID: 20000235}
\textbf{Source problem code:} catalogue entry
\subsection{Definitions and mathematical statement}
This source item records theorem context rather than the following open question. Work over $\C$, with real structures where stated. Let $L=\mathcal O_{\PP^1}(d)$, $0\leq r\leq d$, and
\[
 G=\operatorname{Gr}(r+1,H^0(\PP^1,L)),\qquad N=\dim G=(r+1)(d-r).
\]
A $g^r_d$ is an $(r+1)$-dimensional subspace of sections. At $x\in\PP^1$, the vanishing flag has $j$-dimensional term $F_j(x)=H^0(L(-(d+1-j)x))$. A partition $\lambda$ fitting the $(r+1)\times(d-r)$ rectangle defines the Schubert condition $\Omega_\lambda(F(x))$, of codimension $|\lambda|$. For distinct points $x_1,\ldots,x_s$, form the scheme-theoretic intersection
\[
 S=\bigcap_{i=1}^s\Omega_{\lambda_i}(F(x_i)),\qquad
 m=N-\sum_i|\lambda_i|.
\]
The distinct-point and nonempty-intersection conventions must be made explicit; the extracted text does not fully state them.

The source lists: Theorem 10.1, attributed to Eisenbud--Harris in the 1980s, that a nonempty such intersection has the expected codimension $\sum_i|\lambda_i|$; Theorem 10.2, attributed to Mukhin--Tarasov--Varchenko, that if all $x_i$ are real and $m=0$, then the intersection consists of real reduced points; and Theorem 10.3, that if all $x_i$ are real and $m=1$, then
\[
 S(\R)\subset S_{\mathrm{sm}},
\]
where smoothness is on the complex scheme. The supplied research summary identifies the last assertion as Levinson's theorem. It does not assert that every complex point of this curve is smooth. The separately recorded Problem 10.4 asks that further question and is not part of this item.

The report adds a higher-dimensional residual-Wronski certificate. The Wronskian of a basis $f_0,\ldots,f_r$ is the determinant of their successive derivatives, interpreted as a binary form of degree $N$. Its zero divisor contains $\sum_i|\lambda_i|[x_i]$ on $S$. Removing that fixed divisor gives the reported residual Wronski morphism
\[
 w:S\longrightarrow\operatorname{Sym}^m(\PP^1)\cong\PP^m.
\]
For a nonempty expected-dimensional intersection at distinct real points, the report claims that over every squarefree totally real residual divisor disjoint from those points, the fiber is transverse and all real, $S$ is smooth at every fiber point, and $dw$ is an isomorphism there. Each fixed root-distribution chamber is reported to split into $\delta$ trivial real sheets, where $\delta$ denotes the relevant fiber degree, not explicitly defined in the record.

This certificate does not cover every real point in residual dimension at least two. The report's example is the real plane $\langle z,z^2-1\rangle$ in $\operatorname{Gr}(2,3)$, whose Wronskian is $z^2+1$ and has nonreal roots. It is an obstruction to coverage by this certificate, not a singular-point example or a disproof of real smoothness.
\subsection{Short English statement}
Record expected dimension and real smoothness results for intersections of vanishing-flag Schubert conditions on linear series on $\PP^1$. This is theorem context, and the complex-smoothness question is a separate record. The report adds a residual-Wronski smoothness certificate over distinct real residual roots, while a real Grassmannian example shows that this certificate need not cover all real points.
\subsection{Sources}
\begin{itemize}
\item Local registry: \url{data/unsolvedmath/problems.json}, permanent ID $20000235$, source code \texttt{catalogue entry}; original context, supplied OCR corrections and research summary.
\item AIM workshop ``Degenerations in algebraic geometry'', Problem 10, suggested by ``Jake'': \url{https://aimath.org/pastworkshops/degenalggeomproblems.pdf}. Canonical location \texttt{aim-algebraic-geometry-notes.json}, \texttt{notes[234]}. No PDF was opened here.
\item Source attributions: Eisenbud--Harris, Theorem 10.1; Mukhin--Tarasov--Varchenko, Theorem 10.2 (``Muken'' corrected as supplied); Levinson for Theorem 10.3 according to the summary. No full bibliographic references are supplied; original literature field empty.
\item Supplied verification says the complex-smoothness question is the next record, Problem 10.4. Research attempt 1 is context only, status at run ``not a problem'', registry status open and classification \texttt{OPEN-TRIAGE}; rigor metadata ``proved'', novelty confidence low. The non-peer-reviewed report requires independent verification before citation.
\end{itemize}
\end{document}
